\par\Needspace{4\baselineskip}
\begin{center}
\hypertarget{native-052-article035}{}\pdfbookmark[2]{35.}{article-035}35.
\end{center}
Our \textsc{Bessel} derived analytical expressions for the first three coefficients of the second series, $w^{\prime}$, $w^{\prime\prime}$, $w^{\prime\prime\prime}$, and added a table giving the numerical values of the first two, $w^{\prime}$ and $w^{\prime\prime}$, at each degree of the argument $w$ (\textsc{von Zach} Monatliche Correspondenz, vol. XII. p. 197). A table for the first coefficient $w^{\prime}$ had already been computed by \textsc{Simpson} and is appended to the work of the celebrated \textsc{Olbers} cited above. In most cases this method, with the aid of \textsc{Bessel}'s table, makes it possible to \SourceLeaf{053}determine the true anomaly from the time with sufficient precision. Its remaining shortcomings come down chiefly to the following points.

I. In the inverse problem, namely determining the time from the true anomaly, we have to resort to what amounts to an indirect method and find $w$ from $v$ by trial. To remedy this difficulty, the first series would have to be treated in the same way as the second. It is easy to see that $-v^{\prime}$ is the same function of $v$ as $w^{\prime}$ is of $w$, so that the table for $w^{\prime}$ can also serve for $v^{\prime}$ if the sign alone is changed. A table for $v^{\prime\prime}$ would thus be all that is needed to solve both problems with equal precision.

II. There may indeed be cases where the eccentricity differs so little from unity that the general methods described above appear to give insufficient precision, while the difference is still too large for us safely to neglect the effect of the third and higher powers of $\delta$ in the special method just outlined. In hyperbolic motion particularly, there may be cases in which an error of several arcseconds is unavoidable if we use only ordinary tables giving seven figures, whether we use the general methods or this special method. Though such cases seldom occur in practice, there would surely seem to be something lacking if it took larger tables, which are known to be among the rarer books, to determine the true anomaly to within $0^{\prime\prime}{,}1$ or at least $0^{\prime\prime}{,}2$ in \emph{all} cases. We therefore hope it will be worthwhile to describe a special method we have long used. It also applies even when the eccentricity differs considerably from unity. In this respect at least, it is generally applicable.

\par\Needspace{4\baselineskip}
\begin{center}
\hypertarget{native-053-article036}{}\pdfbookmark[2]{36.}{article-036}36.
\end{center}
Before we begin to explain this method, it is worth noting that the uncertainty of the general methods given above disappears of itself in orbits that approach a parabola as soon as $E$ or $F$ reaches a considerable magnitude. This occurs only at large distances from the Sun. To show this, we write the greatest possible error in an ellipse, which we found in Article 32, IV as $\frac{3\omega ea\sin v}{\lambda r}\cdot206265^{\prime\prime}$, in the form
\[
\frac{3\omega e\BareRoot{(1-ee)}.\sin E}{\lambda(1-e\cos E)^2}\cdot206265^{\prime\prime},
\]
\SourceLeaf{054}from which it is immediately clear that the error is always confined within narrow bounds as soon as $E$ has reached a considerable value, or as soon as $\cos E$ has moved farther from unity, however great the eccentricity may be.

\Needspace{10\baselineskip}
This will be clearer still from the following table, in which we have computed the maximum numerical value of the formula for certain fixed values (for seven decimal figures).
\begin{center}
\begin{tabular}{r|lr}
$E=10^0$ & maximum error $=$ & $3^{\prime\prime}{,}04$\\
$20$ & & $0{,}76$\\
$30$ & & $0{,}34$\\
$40$ & & $0{,}19$\\
$50$ & & $0{,}12$\\
$60$ & & $0{,}08$
\end{tabular}
\end{center}
\noindent The situation is similar for the hyperbola. This is immediately apparent when we write the expression derived in Article 32, VII in the form
\[
\frac{\omega\cos F.(\cos F+3e\sin F)\BareRoot{(ee-1)}}{\lambda(e-\cos F)^2}206265^{\prime\prime}.
\]
The following table gives the maximum values of this expression for certain fixed values of $F$.
\begin{center}
\begin{tabular}{r|r|r|r}
\hline
$F$ & \multicolumn{2}{c|}{$u$} & maximum error\\
\hline
$10^0$ & $1{,}192$ & $0{,}839$ & $8^{\prime\prime}{,}66$\\
$20$ & $1{,}428$ & $0{,}700$ & $1{,}38$\\
$30$ & $1{,}732$ & $0{,}577$ & $0{,}47$\\
$40$ & $2{,}144$ & $0{,}466$ & $0{,}22$\\
$50$ & $2{,}747$ & $0{,}364$ & $0{,}11$\\
$60$ & $3{,}732$ & $0{,}268$ & $0{,}06$\\
$70$ & $5{,}671$ & $0{,}176$ & $0{,}03$
\end{tabular}
\end{center}
\noindent Thus, whenever $E$ or $F$ exceeds $40^0$ or $50^0$ (a case that will seldom occur in orbits that differ little from a parabola, since celestial bodies travelling in such orbits generally escape our view at such great distances from the Sun), there is no reason to abandon the general method. Moreover, in such a case the series we discussed in Article 34 would also converge too slowly. The fact that the method we are about to explain is chiefly suited to cases where $E$ or $F$ has yet to exceed moderate values can therefore in no way be regarded as a defect.
